\chapter*{Abstract}
\addcontentsline{toc}{chapter}{Abstract}
\label{pre:abstract}

\begingroup
\setstretch{1.1}
\setlength{\parindent}{0pt}
\setlength{\parskip}{0.7em}

% Traducción de trabajo del título; revisar junto con el abstract definitivo.
\textbf{Title:} Deep learning-assisted transmission muographic reconstruction through predictive enhancement of low-statistics images for volcanic exploration.\footnote{Undergraduate thesis, research modality.}

\textbf{Authors:} \AutorUno; \AutorDos.\footnote{Universidad Industrial de Santander. Faculty of Physical and Mechanical Engineering. School of Systems Engineering and Informatics. Supervisor: \DirectorNombre, PhD candidate. Co-supervisor: \CodirectorNombre, PhD.}

\textbf{Keywords:} transmission muography, volcanic exploration, deep learning, Monte Carlo simulation, image enhancement, low statistics.

\textbf{Description:}

% Faithful translation of the Spanish description; keep both consistent.
Transmission muography estimates the density distribution inside a volcano from the attenuation of the cosmic-ray muon flux, but the low natural flux requires exposures of weeks or months before density contrasts can be distinguished from counting noise. This work investigated whether a deep learning model trained on synthetic data can transform low-statistics muographic images into images of quality comparable to high-statistics ones and thereby reduce the effective acquisition time in volcanic exploration. Monte Carlo simulations with Geant4 and EcoMug produced 114 references of nineteen idealized volcanic scenarios and, through binomial thinning, 9,120 input--reference pairs with 50, 10, 5 and 1\,\% of the counts, partitioned by geometry. With 5\,\% of the counts, a conditional Wasserstein generative adversarial network (cGAN), a U-Net baseline and six continuations with regional and boundary losses were trained and later evaluated at the four fractions. In a descriptive evaluation on potentially exposed data, the cGAN increased the test SSIM from 0.7643 to 0.8342 and the PSNR from 24.258 to 31.170~dB with respect to its input, below the U-Net, and outperformed on both metrics the unprocessed input with twice the counts in all twenty-four reserved views, but the regional count bias rose from 1.832 to 14.566\,\%. None of the 768 model--fraction--view combinations jointly reached the SSIM and PSNR of the unprocessed input at 50\,\%, so the reduction factor remained undetermined under that criterion. Local indicators, computed on anomalies of known location, showed simultaneous improvements and deteriorations and do not demonstrate the detection of unknown anomalies. It is concluded that the synthetic dataset fulfilled its objective, that comparable quality was only partially supported, and that neither an acquisition-time reduction factor nor quantitative opacity equivalence was demonstrated. The results are limited to synthetic scenarios, without a physical detector or field data.

\endgroup
